%% kapa: exam=1 round=19 year=2008 period=2 form=A subject=economics subject_ko=경제원론 seq=7 number=7
%% source: Q-Net 감정평가사 기출문제 (cst00309) artlSeq=1001962 file=extracted/1008/1001962/2008년 제19회 감정평가사 1차시험 기출문제/2교시.hwp
%% conversion: kordoc 4.24.51 -> tools/gen1.py (auto); review=proofread-sonnet; flags=none
\begin{question}{7}
어떤 개방경제의 국민소득결정 모형이 다음과 같이 주어져 있다. 정부지출을 100에서 200으로 늘리면 균형국민소득은 얼마나 늘어나는가?
\medskip $Y = C + I + G + (EX - IM)$ \par
$C = 200 + 0.6(Y - T),~I = 100,~G = T = 100,~EX = 100,~IM = 100 + 0.1(Y - T)$ \par
\smallskip $Y$: 국민소득, $C$: 소비, $I$: 투자, $G$: 정부지출, $T$: 조세, $EX$: 수출, $IM$: 수입 \par
\begin{choices}
  \item 100
  \item 150
  \item 200
  \item 250
  \item 300
\end{choices}
\end{question}
\answer{3}
